\subsection{Sucesión de Fibonacci}

Recuérdese la definición de la \hyperref[sec:fibonacci_recursiva]{sucesión de Fibonacci}: 
\begin{align*}
    F(0) &= 0 \\
    F(1) &= 1 \\
    F(n) &= F(n-1) + F(n-2) \quad \text{para } n > 1
\end{align*}
El problema consiste en calcular el enésimo elemento de la sucesión. Para abordar el problema mediante programación dinámica, se debe determinar si ostenta las propiedades necesarias y modelar el problema orientado al paradigma.

\subsubsection{Propiedades}

Primero, de entrada se puede ver que una solución general, $F(n)$, en efecto se construye a partir de soluciones para sus subproblemas, $F(n-1)$ y $F(n-2)$. Así pues, este problema cumple con una propiedad análoga a la subestructura óptima (no se dice que presenta subestructura óptima como tal porque no es un problema de optimización).

Segundo, no es difícil ver que el problema presenta subproblemas superpuestos. Para calcular $F(n)$ es necesario calcular $F(n-1)$ y $F(n-2)$; para calcular $F(n-1)$ es necesario a su vez calcular $F(n-2)$ y $F(n-3)$. En tan solo dos pasos, ya se repitió el cálculo de $F(n-2)$. Eso implica que usar programación dinámica si mejorará la eficiencia.

La figura \ref{fig:fibonacci_arbol_ejecucion} muestra el árbol de ejecución para una entrada de \(6\) y hace evidente cuán superpuestos están los subproblemas. La figura muestra que hay dos subárboles con raíz \(4\), lo que indica que el cómputo completo de \(F(4)\) se repite dos veces. Contando llamados, se puede observar que \(F(3)\) se invoca tres veces, \(F(2)\) cinco, \(F(1)\) ocho y \(F(0)\) cinco.
\begin{center}
\setbox0=\hbox{%
\begin{forest}
  for tree={
    draw, rounded corners,
    font=\ttfamily\small,
    inner sep=4pt,
    l sep=1.5em,
    s sep=0.6em,
    edge={-latex},
  }
  [{6}
    [{5}
      [{4}, name=f4a-root
        [{3}
          [{2}
            [{1}, name=f4a-left]
            [{0}]
          ]
          [{1}]
        ]
        [{2}
          [{1}]
          [{0}, name=f4a-right]
        ]
      ]
      [{3}
        [{2}
          [{1}]
          [{0}]
        ]
        [{1}]
      ]
    ]
    [{4}, name=f4b-root
      [{3}
        [{2}
          [{1}, name=f4b-left]
          [{0}]
        ]
        [{1}]
      ]
      [{2}
        [{1}]
        [{0}, name=f4b-right]
      ]
    ]
  ]
  \node[fit=(f4a-root)(f4a-left)(f4a-right), draw, dashed, red, rounded corners, inner sep=4pt] {};
  \node[fit=(f4b-root)(f4b-left)(f4b-right), draw, dashed, red, rounded corners, inner sep=4pt] {};
\end{forest}%
}
\ifdim\wd0>\textwidth
  \resizebox{\textwidth}{!}{\copy0}
\else
  \makebox[\linewidth]{\copy0}
\fi
\captionof{figure}{Árbol de ejecución de \(F(6)\).}
\label{fig:fibonacci_arbol_ejecucion}
\end{center}

\subsubsection{Modelo}

Fibonacci tiene la bondad de que la relación de recurrencia está en el enunciado. A causa de eso, ya están también están definidos los estados: cada estado es simplemente un valor de la entrada, pues al cambiarlo cambia el (sub)problema que se está resolviendo.

A partir de eso, se puede modelar el grafo de dependencias. Para resolver \(F(0)\) o \(F(1)\) no se necesita nada: son casos base, por lo que sus vértices no tienen predecesores. Para resolver \(F(n)\), se necesita \(F(n-1)\) y \(F(n-2)\). Eso aplica para todo \(n > 1\), por lo que por ejemplo para \(F(2)\) se necesita \(F(1)\) y \(F(0)\) y se trazan arcos \(0\longrightarrow 2\) y \(1 \longrightarrow 2\). Se realiza ese proceso para cuantos estados se quiera:

\begin{center}
\setbox0=\hbox{%
\usetikzlibrary{decorations.markings}
\begin{tikzpicture}[
  vertex/.style={circle, draw, minimum size=0.8cm, inner sep=0pt},
  vertexn/.style={circle, draw, minimum size=0.8cm, inner sep=1pt, font=\small},
  hidden/.style={circle, draw=none, minimum size=0.8cm, inner sep=0pt},
  gapmark/.style={postaction={decorate, decoration={markings, mark=at position 0.5 with {\coordinate (gap);}}}},
  -latex,
]
  \foreach \i in {0,...,3} {
    \node[vertex] (f\i) at (1.3*\i, 0) {\i};
  }
  \node[hidden]  (f4)   at (1.3*4, 0) {$\cdots$};
  \node[vertexn] (fnm2) at (1.3*5, 0) {$n{-}2$};
  \node[vertexn] (fnm1) at (1.3*6, 0) {$n{-}1$};
  \node[vertexn] (fn)   at (1.3*7, 0) {$n$};

  \draw (f0) to[bend right=35] (f2);
  \draw (f1) to[bend right=20] (f2);
  \draw[-] (f2) to[bend right=35] (f4);
  \draw[-] (f3) to[bend right=20] (f4);
  \draw (f4) to[bend right=35] (fnm1);
  \draw (fnm2) to[bend right=20] (fnm1);

  \draw (f1) to[bend left=35] (f3);
  \draw (f2) to[bend left=20] (f3);
  \draw[gapmark] (f3) to[bend left=35] (fnm2);
  \fill[white] (gap) circle (0.40);
  \draw (f4) to[bend left=20] (fnm2);
  \draw (fnm2) to[bend left=35] (fn);
  \draw (fnm1) to[bend left=20] (fn);
\end{tikzpicture}%
}
\ifdim\wd0>\textwidth
  \resizebox{\textwidth}{!}{\copy0}
\else
  \makebox[\linewidth]{\copy0}
\fi
\captionof{figure}{Grafo de dependencias de \(F(n)\).}
\label{fig:fibonacci_grafo_dependencias}
\end{center}

Con eso, el problema está modelado por completo. De este punto en adelante, implementar los algoritmos de programación dinámica es simplemente plasmar los modelos.

\subsubsection{Solución top-down}
\label{sec:fibonacci_top_down}

Se crea un arreglo donde se memoizarán las respuestas. La enésima posición corresponde al cálculo del enésimo número de Fibonacci. Inicialmente se llena con $-1$ para indicar que no se ha computado ningún número.

La función principal pasa el arreglo a la función auxiliar, que tiene la \hyperref[impl:fibonacci_recursivo]{implementación recursiva ingenua} con dos modificaciones: se memoiza el cómputo si aún no está en el arreglo y siempre se retorna la solución memoizada.
\begin{codeblock}[Fibonacci top-down.][impl:fibonacci_top_down]
\begin{pseudocode}
def fibonacci_top_down(n: int) -> int:
    memo = |\textrm{arreglo de enteros de tamaño }$n+1$|
    |\br{cbi}|memo[0] = 0
    memo[1] = 1|\br{cbf}|

    |\br{lli}|for i=2 to n:
        memo[i] = -1|\br{llf}|

    return fibonacci_memo(n, |\bxl{auxf}|memo)|\bxr{auxf}|

def fibonacci_memo(n: int, memo: arreglo[int]) -> int:
    |\bxl{memo}|if memo[n] < 0:|\bxr{memo}|
        memo[n] = (
            fibonacci_memo(n - 1, memo) + 
            fibonacci_memo(n - 2, memo)
        )
    |\bxl{rmemo}|return memo[n]|\bxr{rmemo}|
\end{pseudocode}
\begin{annotations}
    \codebrace[width=12cm]{cbi}{cbf}{Se añaden los casos base directamente al arreglo, pues son soluciones ya computadas.}
    \codebrace[width=9cm]{lli}{llf}{Se llena el resto del arreglo con -1 para indicar que esas soluciones no se han calculado. Es inambiguo porque los términos calculados serán positivos.}
    \codenote[box=true, width=4.5cm]{auxf}{La función auxiliar recursiva recibe siempre la misma estructura y la va llenando.}
    \codenote[width=10cm]{memo}{La estructura aún no contiene la solución para este valor, entonces, se memoiza.}
    \codenote[width=6.5cm]{rmemo}{Se retorna siempre el valor memoizado (que puede que se haya calculado en ese llamado recursivo o en uno anterior).}
\end{annotations}
\end{codeblock}

La memoización hace que cada llamado recursivo se repita como máximo dos veces: la primera vez (en orden de ejecución, es decir, en orden de un recorrido DFS sobre el árbol) se realiza el cómputo, por lo que se genera un subárbol; pero, la segunda vez, se retorna enseguida. Eso genera el siguiente árbol de ejecución:

\begin{center}
\begin{forest}
  for tree={
    draw, rounded corners,
    font=\ttfamily\small,
    inner sep=4pt,
    l sep=1.5em,
    s sep=1.2em,
    edge={-latex},
  }
  [{6}
    [{5}
      [{4}
        [{3}
          [{2}
            [{1}]
            [{0}]
          ]
          [{1}]
        ]
        [{2}]
      ]
      [{3}]
    ]
    [{4}]
  ]
\end{forest}
\captionof{figure}{Árbol de ejecución de \(F(6)\) con memoización top-down.}
\label{fig:fibonacci_arbol_top_down}
\end{center}

Nótese que la cantidad de llamados recursivos es lineal con respecto a la entrada, hay máximo dos llamados por cada elemento de la sucesión de Fibonacci. Para cualquier entrada \(n > 1\), el árbol tiene \(2n - 1\) vértices. Por ende el costo recursivo es lineal, sumado al costo no recursivo que es constante (solo accesos a un arreglo), la complejidad temporal de la solución es \(\Theta(n)\).

Respecto a la complejidad espacial, el arreglo en el que se memoizan las soluciones es del tamaño de la entrada, por lo cual su costo es lineal. Aunado a eso, la altura del árbol no cambia, por lo que la profundidad máxima de la pila de llamados es la misma que la de la solución recursiva ingenua, también lineal. Eso hace que la complejidad espacial sea \(\Theta(n)\).

\subsubsection{Solución bottom-up}
\label{sec:fibonacci_bottom_up}

El método bottom-up resuelve primero los problemas más pequeños. En este caso, eso traduce naturalmente a calcular los términos de la sucesión en orden ascendente, siguiendo el orden topológico del \hyperref[fig:fibonacci_grafo_dependencias]{grafo de dependencias}.

\begin{codeblock}[Fibonacci bottom-up.][impl:fibonacci_bottom_up]
\begin{pseudocode}
def fibonacci(n: int) -> int:
    tab = |\textrm{arreglo de }|int|\textrm{ de tamaño }$n+1$|
    |\br{nds}|tab[0] = 0
    tab[1] = 1|\br{ndf}|

    |\br{grs}|for i=2 to n:
        tab[i] = tab[i-1] + tab[i-2]|\br{grf}|

    return tab[n]
\end{pseudocode}
\begin{annotations}
    \codebrace[width=14cm]{nds}{ndf}{Solo las soluciones para \(0\) y \(1\) se añaden directamente a la tabla porque no tienen dependencias, de acuerdo con el grafo.}
    \codebrace[width=8cm]{grs}{grf}{Se utilizan las dependencias del grafo, en orden, para llenar la tabla.}
\end{annotations}
\end{codeblock}

Se obtiene un algoritmo con un solo ciclo. Su complejidad temporal, a ojímetro, es $\Theta(n)$. Su complejidad espacial es solo lo que ocupa la tabulación, que es un arreglo del tamaño de la entrada: $\Theta(n)$ también.

Como para calcular el enésimo número de Fibonacci es necesario calcular todos los anteriores, el algoritmo bottom-up es (ligeramente) más eficiente que el algoritmo top-down: ambos calculan todos los términos previos, pero el bottom-up evita el \hyperref[sec:sobrecosto_recursion]{sobrecosto de la recursión}.

\subsubsection{Optimización en espacio}
\label{sec:fibonacci_bottom_up_optimizada}

Si se observa cualquier vértice (mayor a \(1\)) del \hyperref[fig:fibonacci_grafo_dependencias]{grafo de dependencias}, resulta evidente que solo se necesitan las últimas dos soluciones para computar la próxima. La tabulación se puede realizar entonces empleando solo dos variables, en lugar de un arreglo. 

Eso hace que la complejidad espacial no dependa del tamaño de la entrada, reduciéndola a $\Theta(1)$.

\begin{codeblock}[Fibonacci bottom-up optimizado en espacio.][impl:fibonacci_bottom_up_espacio]
\begin{pseudocode}
def fibonacci(n: int) -> int:
    |\bxl{sol}|solution = n|\bxr{sol}|
    prev_prev_sol = 0
    prev_sol = 1

    for i=2 to n:
        solution = prev_sol + prev_prev_sol
        prev_prev_sol = prev_sol
        prev_sol = solution

    return solution
\end{pseudocode}
\begin{annotations}
    \codenote[width=14cm]{sol}{Se inicializa la solución como \(n\) para que la función retorne eso si \(n \leq 1\), sin entrar al ciclo.}
\end{annotations}
\end{codeblock}

\practicar[leetcode]{Fibonacci Number}{https://leetcode.com/problems/fibonacci-number/}
